Antibubbles -- liquid droplets encapsulated by a thin gas film within a continuous liquid phase -- hold promise for ultrasound imaging, drug delivery, and microgravity fluid management, yet no existing method generates them with both high efficiency and repeatability. Here we overcome this bottleneck by replacing the random foam layer of the traditional foam-based method with a 3D-printed triangular-prism frame that fixes a single Plateau border (PB) in space, transforming a stochastic process into a controlled physical system. Across 199 experiments spanning droplet radii R = 1.0-2.2 mm, impact velocities v0 = 0.1-1.8 m/s, and ambient pressures P = 0.1-1.0 atm, we identify four universal stages of droplet-PB interaction and three outcomes: packing (successful encapsulation), outer coalescence, and inner coalescence. From energy conservation and gas-film lubrication theory, we derive two dimensionless criteria -- a We-Bo relation [We + 2 Bo = 0.038 +/- 0.004] and a characteristic-time ratio t_p/t_d < 1 -- that together predict 97% of experimental outcomes. The We-Bo criterion quantifies why the inclined PB geometry is twice as efficient as horizontal-interface methods, while the t/t criterion reveals a pressure-dependent failure mode relevant to reduced-pressure environments. This work establishes a predictive framework for antibubble formation and provides a deterministic experimental platform for studying confined transport in foam microchannels.